%==============================================================================
%  Project brief for Prof. Adithya Vadapalli
%  Trusted delegated payments for autonomous agents
%
%  Written for a cryptographer who does not work on payments.
%  Section 2 is the payments background. Everything after it assumes only that.
%
%  Compile: pdflatex project-brief-crypto.tex  (twice)
%==============================================================================
\documentclass[journal,onecolumn,11pt]{IEEEtran}

\usepackage[utf8]{inputenc}
\usepackage[T1]{fontenc}
\usepackage[a4paper,margin=1in]{geometry}
\usepackage{booktabs}
\usepackage{enumitem}
\usepackage{xcolor}
\usepackage{amsmath,amssymb}
\usepackage{url}
\usepackage[hidelinks]{hyperref}

\newcommand{\rs}{Rs.\,}
\definecolor{accent}{HTML}{0D5C66}
\definecolor{warnc}{HTML}{8A5A10}
\definecolor{linkc}{HTML}{15505C}
\hypersetup{colorlinks=true,urlcolor=linkc,linkcolor=linkc,citecolor=linkc}

\newenvironment{keybox}
  {\begin{quote}\color{accent}\rule{\linewidth}{0.8pt}\vspace{-0.4em}\begin{itshape}}
  {\end{itshape}\vspace{-0.2em}\rule{\linewidth}{0.8pt}\end{quote}}

\newcommand{\q}[1]{\textcolor{warnc}{\textbf{Open question.}} #1}

\title{Trusted Delegated Payments for Autonomous Agents\\
\large A project brief, written for a cryptographer}

\author{Ashutosh Anand\\
\textit{Indian Institute of Technology Kanpur}\\
Supervisor: Prof.\ Vimal Kumar \quad$\cdot$\quad 21 August 2026}

\begin{document}
\maketitle

\begin{abstract}
\noindent I am writing this because you asked me a question I could not answer properly
when we spoke. You said a credential can be stored locally. Then you asked how that lets
me verify the agent. How do I tell a good agent from a bad one. How do I call it trusted.

Section~\ref{sec:verify} is my honest answer. It turned out to be more interesting than
I expected because part of the answer is negative. Sections~\ref{sec:bg}
and~\ref{sec:proto} are the payments background you need before that answer means
anything. I assume you know no payments. I have described each mechanism by what it
actually guarantees rather than by what its vendor says about it.
Section~\ref{sec:crypto} is where I think your work and this problem meet. That is the
section I most want your reaction to.
\end{abstract}

\begin{IEEEkeywords}
Delegated authorisation, agent attestation, trusted execution, secure computation,
private information retrieval, payment mandates, UPI, revocation.
\end{IEEEkeywords}

\section{What I am building}

\IEEEPARstart{A}{n} AI agent sits inside a device. It could be an appliance or a vehicle
charger or a diagnostic instrument. The agent notices it needs something and buys it.
Nobody is watching. The owner approved a spending arrangement weeks ago and is not
present when the payment happens.

I am building the layer that decides whether that payment should go through. It also
produces the evidence needed afterwards if it should not have.

\begin{keybox}
Every payment rail shipping for agents today enforces how much an agent may spend. None
of them verifies what it actually bought.
\end{keybox}

Here is what that means concretely. A mandate authorises \rs 2{,}000 per purchase. The
agent asks to pay \rs 600. The owner approved a basket worth \rs 118. No limit is
violated. No rule is broken. The owner loses \rs 482 and the dispute system has no
category to file it under.

I have this running in a simulator. It takes two commands to show.

\subsection{Scope}
My work starts when the basket is finished and the agent asks to pay. It ends when the
money has moved and the evidence is recorded. Product discovery and price comparison and
supplier choice all happen before that. They are out of scope.

I restrict the problem to repetitive purchases. The same consumable is reordered again
and again. This is what makes advance authorisation meaningful. The owner cannot approve
each purchase but they can approve the pattern.

\section{Payments background}\label{sec:bg}

This section assumes you know nothing about payments. Skip whatever is familiar.

\subsection{UPI}
UPI is India's real time retail payment system. It is operated by
\href{https://www.npci.org.in/what-we-do/upi}{NPCI} which is a not for profit consortium
of banks. A user is addressed by a virtual payment address like \texttt{name@bank}. That
address resolves to a bank account. A payment is a message routed through a central
switch. The switch debits one bank and credits another. It settles in seconds.

Two properties matter here. UPI is zero MDR for most merchant transactions. There is
no interchange fee so the usual payments revenue model does not exist here.
Authentication is a four or six digit UPI PIN entered on a device bound to the user. The
security model assumes a human is present.

\subsection{Mandates}
A mandate is standing authorisation to debit an account without fresh authentication
each time. Two kinds matter.

UPI Autopay is a recurring mandate. It is bound to one payee and one amount ceiling and
a schedule. It was built for subscriptions. It cannot express ``pay whichever supplier
turns out to be cheapest'' because the payee is fixed when the mandate is created.

UPI Reserve Pay is formally called Single Block Multiple Debits. It blocks funds inside
the user's own account. A merchant then debits against that blocked pool repeatedly
until it is empty or expires. The ceiling is \rs 10{,}000 and validity is 90 days. There
can be one active block per merchant per customer. That last constraint shapes my whole
architecture.

Here is the fact that matters for security. Authentication happens once when the block
is created. Every debit after that is unauthenticated by the human. The rail carries no
field saying who or what started a debit. So it cannot tell an agent from a person.

\subsection{Cards and why they are different}
Card networks spent two decades building infrastructure that assumes an adversary. UPI
never needed that. Cards have tokenisation and network side authorisation policy and a
formal dispute process with defined liability shifts and a 120 day chargeback window.

Tokenisation replaces a card number with a network issued token scoped to a context.
This is what Apple Pay does. The token is a reference and not a bearer instrument. It
holds no value. It de-tokenises inside the network back to the real account.

\subsection{The asymmetry that shapes everything}
UPI is cheap and everywhere in India and has no notion of an agent. Cards are expensive
but they carry an agent identifier and enforce scope at the network and have a dispute
process.

So my system spans both rails. The owner picks which one funds a given agent.

\section{The protocol landscape}\label{sec:proto}

Four systems shipped in the last eighteen months. I describe each one by what it
actually guarantees cryptographically.

\subsection{AP2 from Google, September 2025}
\href{https://ap2-protocol.org}{AP2} defines three signed artifacts as
\href{https://www.w3.org/TR/vc-data-model-2.0/}{W3C Verifiable Credentials}. An Intent
Mandate says what the user asked for. A Cart Mandate says what the agent selected. A
Payment Mandate authorises the movement. Each artifact is non repudiable.

AP2 defines the envelope. It does not define a constraint language for what goes inside
the envelope. So the question ``does this cart satisfy this intent'' is left unspecified.
That gap is where my work sits.

\subsection{Mastercard Agent Pay and Verifiable Intent}
\href{https://www.mastercard.com/us/en/business/artificial-intelligence/mastercard-agent-pay.html}{Agentic
Tokens} extend card tokenisation. They add an agent identifier and a scope object holding
spend caps and merchant categories and expiry. The network validates that scope at every
single authorisation.

Verifiable Intent was open sourced in March 2026. It is a credential that travels with
the authorisation and cryptographically links the transaction to the human mandate behind
it.

Mastercard calls their framing Know Your Agent. It asks three questions and they map
directly onto yours. Is the controlling human verified. Are the agent's keys in TEE
hardware. Is the delegation scope well defined.

\subsection{P3P from Pine Labs, June 2026}
\href{https://www.pinelabs.com/docs/online-payments/ai/p3p}{P3P} is the closest system to
my project and it is live in production. It runs agentic payments over UPI Reserve Pay
using HTTP 402. An identity and delegation layer called Grantex supplies agent identity
and spend limits and an audit trail.

\q{Their public documentation stops at a quickstart. It does not describe any binding
between a payment and an approved order. If such a binding exists then a large part of my
contribution is already taken. This is my highest priority open item and I have not yet
resolved it.}

\subsection{x402 and MPP}
Both revive HTTP status code 402. \href{https://github.com/coinbase/x402}{x402} from
Coinbase settles USDC on chain using EIP-3009 gasless authorisations verified by a
facilitator. \href{https://stripe.com/blog/machine-payments-protocol}{MPP} from Stripe
and Tempo is a general HTTP authentication scheme on an IETF track with payment methods
defined as separate specifications.

Both solve the problem of paying for an HTTP resource. There the thing purchased is the
request itself. Neither needs an order record because there is no basket. They are
adjacent to my problem and not competing with it.

\section{My answer to your question}\label{sec:verify}

You said a credential can be stored locally. You asked how that lets me verify the agent
and separate it from a bad agent and treat it as trusted.

A local credential gives me one of the four properties I need. The fourth one cannot be
obtained by cryptography at all. Separating the four is most of the answer.

\subsection{Four properties that get confused}

\paragraph{One. Authentication. Which key signed this?}
The agent holds an Ed25519 keypair and signs each request. This establishes that
something holding that key produced the message.

This is what a locally stored credential gives me. It is necessary. It is also the
weakest of the four. Against an adversary with host access it establishes nothing at all
because a software key can be read out of memory and used elsewhere. An agent
authenticated this way looks identical to an attacker who copied its key.

\paragraph{Two. Binding. Is the key where I think it is?}
Generate and hold the key inside hardware isolation. That means a TEE like ARM TrustZone
or Intel TDX or AMD SEV-SNP. Or a Secure Enclave. Or a TPM. The private key never leaves.
The environment signs on request.

Now the claim is stronger. Before it was that something has the key. Now it is that
something on this specific device has the key. Extracting the key means defeating the
hardware. This is exactly what Mastercard's Know Your Agent framework is asking about
when it asks whether agent keys sit in TEE hardware.

\paragraph{Three. Attestation. What code is running?}
This is the interesting one. A TEE can produce a signed quote. The quote contains a
measurement which is a hash of the code loaded into the enclave. It is signed by a key
rooted in the hardware manufacturer's PKI. A remote verifier checks the signature chain
and compares the measurement against a known good value.

This gives me what authentication cannot. It is evidence that this specific software
build is running unmodified inside genuine hardware. Attestation is the mechanism that
separates a legitimate agent from a modified one.

It needs a registry. A measurement only means something if somebody publishes that
measurement $m$ corresponds to approved agent build $B$. NPCI's proposed Unified Agent
Protocol and Mastercard's Know Your Agent are both reaching for this. Neither has
shipped a specification.

\paragraph{Four. Correctness. Did the agent do the right thing?}
Here the stack stops helping me. This is the part I find genuinely interesting.

\begin{keybox}
Attestation of an AI agent certifies strictly less than attestation of a deterministic
program. The same attested binary produces different behaviour depending on input that
the attestation does not cover.
\end{keybox}

Picture an agent that is attested and unmodified and correctly keyed. It reads a
supplier's product description. The description contains adversarial text. The agent buys
the wrong thing.

Every cryptographic property above still holds. The measurement is correct. The signature
verifies. The hardware is genuine. The purchase is wrong.

Attestation covers the code and at best the model weights. It cannot cover the inference
because the inference depends on runtime input from an untrusted source. For an ordinary
program I would call that input validation and move it inside the measured boundary. For
an agent whose entire job is reading untrusted text and acting on it that move is not
available to me.

\subsection{What follows for the architecture}
Property four is unreachable by attestation. So I have to get it structurally instead. I
constrain what the agent is allowed to do so that being wrong is not enough to cause a
loss.

This is the central mechanism of the project. The owner signs a constraint in advance.
Before settlement a deterministic check tests whether the requested payment satisfies
that constraint. A compromised agent can still request anything. It cannot get settlement
for a request that fails the check.

Two design rules follow. I got both of them wrong on my first attempt.

\begin{enumerate}[leftmargin=*]
\item Enforcement must sit outside the agent's trust boundary. If the agent checks its
own limit then an injected agent moves its own limit. The check runs in a party the agent
cannot influence.
\item The check must be deterministic. No model sits in the approval path. An issuer
cannot decline a payment on a probability. It needs a reason that can be written into a
dispute file. This also means the check itself is not vulnerable to the same class of
attack as the agent.
\end{enumerate}

\subsection{A result that surprised me}
I implemented the order as a signed artifact. Then I tested three signing arrangements.
Signed by the agent. Signed by the merchant. Signed by both.

All three approve a self consistent inflated order. If the line items genuinely add up to
\rs 600 and the signature verifies then nothing objects. The order is the only surviving
record of what was agreed so there is nothing left to compare it against. What actually
refuses it is a per payment ceiling close to the expected basket size.

The lesson generalises and it is your question one layer up. A signature establishes
authorship. It does not establish truth. I have found this a useful discipline for
thinking about the rest of the design.

\section{Where I think your work meets this problem}\label{sec:crypto}

Everything up to here is mostly systems engineering. This section is what I think is
actually research. All of it sits closer to your area than to mine.

I have read enough of your work to know roughly where you sit. Your
\href{https://avadapal.github.io/}{page} lists private information retrieval and secure
multiparty computation and oblivious RAM and metadata protection. Your thesis was on
distributed point functions. Three of the four problems below touch that directly and I
would not have found them without reading it.

\subsection{Problem A. The evidence trail is a surveillance database}
To settle disputes my system records what every machine bought. It binds each purchase to
an owner identity and keeps it indefinitely. That is a complete record of a household's
or a firm's consumption. It is held by an operator with a commercial interest in it.

For a dispute I need to reveal exactly one purchase to exactly one arbiter. Everything
else should stay hidden from the arbiter and from the merchant and ideally from the
operator holding the log.

\q{Can I build the dispute evidence so the owner proves a specific payment violated a
specific mandate without revealing the mandate's other terms and without revealing any
other purchase and without the operator holding a readable history?}

This looks to me like selective disclosure over a commitment. Maybe a Merkle committed
log with zero knowledge membership and predicate proofs. Maybe anonymous credential
techniques applied to the mandate. I do not know which is right. I also do not know
whether verification cost fits inside an authorisation path that must finish well under a
second.

\subsection{Problem B. Checking a mandate without learning it}
Right now my operator sees both the mandate and the order because it has to evaluate one
against the other. That is a lot of trust placed in a party that has no business knowing
what a household buys.

The check itself is narrow and arithmetic. A comparison against an amount ceiling.
Membership in a merchant set. Membership in a category set. A period budget subtraction.
A validity window. These are exactly the primitives secure computation handles well.

\q{Can the order check run as a two party computation between the agent's host and the
verifier so the verifier learns only the boolean verdict and neither party learns the
other's input?}

The merchant allow list is a private set membership query. That is where I suspect your
distributed point function work applies directly. The realistic obstacle is latency. An
authorisation decision has roughly 150 milliseconds end to end. I do not know what is
achievable in that budget for a set of realistic size and I would rather you told me it
is hopeless now than in a year.

\subsection{Problem C. Revocation}
An owner revokes an agent. That revocation has to reach two rails and any funds already
in flight. A verifier also has to establish at authorisation time that a credential is
still valid.

Naive approaches leak. A status query reveals which agent is transacting and when. This
is the classic tension between revocation freshness and verification cost and privacy.
Accumulators and short lived credentials attack it from different directions.

\q{What revocation construction gives sub second freshness without the status check
itself becoming a tracking channel?}

There is a timing question underneath this with real money attached. If the owner revokes
at time $t$ and a payment authorises at $t + \epsilon$ then who takes the loss. My system
currently emits a signed timestamped revocation record so the question is at least
decidable. The right construction is not obvious to me.

\subsection{Problem D. Metadata}
Even with contents hidden the pattern leaks. Purchase frequency and timing and amount
distribution and the identity of the settling merchant. For a household that reveals
occupancy and travel and health.

This is your listed interest in metadata protection applied to a setting where the
observer is a regulated financial intermediary. I cannot simply remove that observer
because the regulation requires it to exist.

\subsection{An honest split}
The order check and the two rail architecture and the reconciliation logic and the control
surface are engineering. They are substantial and they work. They are not research and I
will not claim they are.

Problems A through D are research. I cannot do them alone. That is why I am asking you.

\section{Threat model}

\subsection{What I assume the adversary can do}
\begin{itemize}[leftmargin=*]
\item The agent can be compromised at runtime by adversarial input. I assume this. I do
not defend against it. Prompt injection is unsolved and I do not claim to prevent it.
\item The merchant is adversarial. It controls all product data reaching the agent and it
has a direct financial incentive to influence what the agent selects.
\item The operator is honest but curious. It follows the protocol and would like to learn
what everyone buys. Problems A and B exist to reduce what it can learn.
\item Rails are trusted for correctness and not for privacy.
\end{itemize}

\subsection{Security goals}
\begin{enumerate}[leftmargin=*]
\item A compromised agent cannot obtain settlement for a payment that violates the
owner's signed constraint.
\item Every settled payment produces evidence sufficient to adjudicate a dispute.
\item Revocation is effective across both rails and any funds in flight.
\item No party learns more about the owner's purchasing than its role requires. This is a
target and I have not achieved it.
\end{enumerate}

Goal four is the research contribution. Goals one through three are built and tested.

\section{What I have already built}

A simulator with UPI shaped rails. It has a double entry ledger where every posting must
sum to zero. A two leg payment with a suspense account between the legs. Injectable
failures. A reconciliation sweep. Ed25519 signed authority and order records. The
deterministic order check. Cross rail virtual card settlement. Revocation. An owner
console. There are 69 tests and none of them need a network.

Four results came out of building it. Each one is asserted by a test rather than argued
in prose.

\begin{enumerate}[leftmargin=*]
\item Signatures on an order do not prevent overcharging under any signer arrangement.
\item Routing funds through an agent controlled card destroys merchant scoping on the
funding leg because the mandate can then only name the card.
\item A failed and refunded purchase permanently consumes the owner's periodic budget
unless headroom is explicitly released on reversal.
\item Revocation is three operations. Most implementations perform one.
\end{enumerate}

\section{What I would like from you}

Not answers. You already told me you would not have them and the interesting parts do not
have known answers yet.

What would help me most is a judgement on whether Problems A through D are stated well. I
might have written down tractable looking versions of problems that are actually trivial
or actually hopeless. I cannot tell that from inside the problem.

After that I would rather work on whichever of the four you find most interesting. A
supervisor who is engaged with the problem is worth more to me than a problem I picked
alone.

\section*{Reading, if it helps}

\textbf{Payments, in the order I found useful}
\begin{itemize}[leftmargin=*,itemsep=1pt]
\item NPCI Operating Circular 228 on Reserve Pay. This is the primitive everything sits
on. \href{https://www.npci.org.in/what-we-do/upi}{npci.org.in}
\item Google. \href{https://ap2-protocol.org}{Agent Payments Protocol}
\item Mastercard.
\href{https://www.mastercard.com/us/en/business/artificial-intelligence/mastercard-agent-pay.html}{Agent
Pay and the agentic token framework}. The Know Your Agent questions in
Section~\ref{sec:verify} come from here.
\item Pine Labs.
\href{https://www.pinelabs.com/docs/online-payments/ai/p3p}{P3P documentation}
\end{itemize}

\textbf{Closer to your area}
\begin{itemize}[leftmargin=*,itemsep=1pt]
\item \href{https://www.w3.org/TR/vc-data-model-2.0/}{W3C Verifiable Credentials Data
Model}. This is the artifact format AP2 uses.
\item Remote attestation literature for Intel TDX and AMD SEV-SNP. This is the basis for
property three in Section~\ref{sec:verify}.
\item \href{https://arxiv.org/abs/2604.15367}{arXiv:2604.15367} on unverified instruction
origin after mandate creation. This establishes that the observation is already known. So
only a mechanism counts as a contribution now.
\end{itemize}

\vspace{1em}
\noindent\rule{\linewidth}{0.4pt}\\
\small The code and full documentation are available. The simulator runs offline with one
dependency. I am happy to walk through any part of it. I would much rather be told early
that a direction is wrong than find out at the end.

\end{document}
